\section{Proofs}\label{app:proofs}

\begin{proof}[Proof of Lemma~\ref{lem:cara-demand}]
The agent with posterior $\mu$ and wealth $W$ facing price $p$ chooses
$x$ to maximise
\[
   V(x) \;=\; \mu\,U\bigl(W+x(1-p)\bigr) + (1-\mu)\,U\bigl(W-xp\bigr).
\]
Setting $V'(x)=0$:
\begin{equation}
   \mu(1-p)\,U'\bigl(W+x(1-p)\bigr) \;=\; (1-\mu)\,p\,U'\bigl(W-xp\bigr).
\label{eq:app-foc}
\end{equation}
Under CARA, $U(W)=-e^{-\alpha W}/\alpha$ and $U'(W)=e^{-\alpha W}$.
Substituting:
\[
   \mu(1-p)\,e^{-\alpha[W+x(1-p)]} \;=\; (1-\mu)\,p\,e^{-\alpha[W-xp]}.
\]
The factor $e^{-\alpha W}$ appears on both sides and cancels (this
cancellation is the defining property of CARA---no wealth effects):
\[
   \mu(1-p)\,e^{-\alpha x(1-p)} \;=\; (1-\mu)\,p\,e^{\alpha xp}.
\]
Dividing both sides by $(1-\mu)\,p$ and collecting the exponentials on the
right:
\[
   \frac{\mu(1-p)}{(1-\mu)\,p}
   \;=\; e^{\alpha x(1-p)+\alpha xp}
   \;=\; e^{\alpha x}.
\]
The left-hand side equals $e^{\logit\mu-\logit p}$ by definition of the
logit function $\logit(q)=\ln\tfrac{q}{1-q}$. Taking logarithms:
$\alpha x = \logit\mu-\logit p$, giving
$x = [\logit\mu-\logit p]/\alpha$.
\end{proof}

\begin{proof}[Proof of Lemma~\ref{lem:crra-demand}]
The first-order condition \eqref{eq:app-foc} with CRRA utility
$U(W)=W^{1-\gamma}/(1-\gamma)$ and marginal utility $U'(W)=W^{-\gamma}$
becomes
\[
   \mu(1-p)\,\bigl(W+(1-p)x\bigr)^{-\gamma}
   \;=\; (1-\mu)\,p\,\bigl(W-px\bigr)^{-\gamma}.
\]
Rearranging the ratio of marginal utilities:
\[
   \left(\frac{W+(1-p)x}{W-px}\right)^{-\gamma}
   \;=\; \frac{(1-\mu)\,p}{\mu\,(1-p)}
   \;=\; e^{\logit p - \logit\mu}.
\]
Raising both sides to the power $-1/\gamma$:
\[
   \frac{W+(1-p)x}{W-px}
   \;=\; \exp\!\left(\frac{\logit\mu-\logit p}{\gamma}\right) \;\equiv\; R.
\]
Cross-multiplying: $W+(1-p)x = R\,(W-px)$, so
$W + (1-p)x = RW - Rpx$, giving
$x\bigl[(1-p)+Rp\bigr] = W(R-1)$, which yields
\eqref{eq:crra-demand}.
\end{proof}

\begin{proof}[Proof of Lemma~\ref{lem:cara-limit}]
Let $z\equiv\logit(\mu_{k})-\logit(p)$ and $\delta\equiv z/\gamma$, so
$R_{k}=e^{\delta}$. For each fixed $(\mu_{k},p)\in(0,1)^{2}$, the
quantity $z$ is fixed and $\delta\to 0$ as $\gamma\to\infty$.
Taylor-expanding the exponential:
\[
   R_{k}-1 \;=\; e^{\delta}-1
   \;=\; \delta+\tfrac{\delta^{2}}{2}+O(\delta^{3}).
\]
Similarly, $(1-p)+R_{k}\,p = (1-p)+p\,e^{\delta}
= 1 + p\,\delta + O(\delta^{2})$. Hence the CRRA demand becomes
\[
   x_{k}^{\mathrm{CRRA}}
   \;=\; \frac{W_{k}\,(R_{k}-1)}{(1-p)+R_{k}\,p}
   \;=\; \frac{W_{k}\,[\delta+O(\delta^{2})]}{1+p\,\delta+O(\delta^{2})}
   \;=\; W_{k}\,\delta\bigl[1+O(\delta)\bigr].
\]
Multiplying both sides by $\gamma$ and substituting
$\gamma\delta=z=\logit(\mu_{k})-\logit(p)$:
\[
   \gamma\,x_{k}^{\mathrm{CRRA}}
   \;=\; W_{k}\,[\logit(\mu_{k})-\logit(p)]\,\bigl[1+O(1/\gamma)\bigr]
   \;\xrightarrow{\gamma\to\infty}\;
   W_{k}\,[\logit(\mu_{k})-\logit(p)],
\]
which is the CARA demand \eqref{eq:cara-demand} with risk-aversion
parameter $\alpha_{k}=\gamma/W_{k}$. The convergence is uniform on
compact subsets of $(0,1)^{2}$ because $\delta=z/\gamma\to 0$ uniformly
when $z$ is bounded.
\end{proof}

\begin{proof}[Proof of Theorem~\ref{thm:cara-unique}]
The first-order condition for the binary-asset problem is
\[
   \mu(1-p)\,U'\!\bigl(W+(1-p)x\bigr) \;=\; (1-\mu)\,p\,U'\!\bigl(W-px\bigr).
\]
Taking logarithms and using $\logit(q)=\ln\tfrac{q}{1-q}$, this
becomes
\begin{equation}
   \ln U'\!\bigl(W+(1-p)x\bigr) - \ln U'\!\bigl(W-px\bigr)
   \;=\; \logit p - \logit\mu
   \;=\; -z,
\label{eq:thm-foc-v2}
\end{equation}
where $z\equiv\logit\mu-\logit p$.

Define $h(W)\equiv\ln U'(W)$, which is well defined and three times
continuously differentiable on $\mathcal{D}$ since $U'>0$.
Now suppose the optimal demand admits the representation
\eqref{eq:linear-demand-rep}, i.e.\ $x^{\star}=z/\alpha(p,W)$ for some
positive function $\alpha$.  Substituting into \eqref{eq:thm-foc-v2}:
\begin{equation}
   h\!\bigl(W+(1-p)z/\alpha\bigr) - h\!\bigl(W-pz/\alpha\bigr) \;=\; -z,
\label{eq:thm-h-v2}
\end{equation}
which must hold for all admissible $(W,p,z)$ with $z\in\R$ ranging freely
as $\mu$ varies in $(0,1)$ at fixed $p$.

\emph{Step 1: First derivative in $z$.} Differentiate
\eqref{eq:thm-h-v2} with respect to $z$ at $z=0$ (where both arguments of
$h$ equal $W$). By the chain rule:
\begin{align*}
   &h'(W)\cdot\frac{1-p}{\alpha(p,W)}
   \;+\; h'(W)\cdot\frac{p}{\alpha(p,W)} \\
   &\quad=\; \frac{h'(W)}{\alpha(p,W)}
   \;\stackrel{!}{=}\; -1.
\end{align*}
(The second term gets a factor $p/\alpha$ rather than $-p/\alpha$
because the derivative of $W-pz/\alpha$ with respect to $z$ is
$-p/\alpha$, and the overall minus sign in front of the second $h$-term
in \eqref{eq:thm-h-v2} produces another sign flip.)
Hence $\alpha(p,W)=-h'(W)$ for all $(p,W)$.
Since $\alpha>0$ by assumption, $h'(W)<0$ for all $W\in\mathcal{D}$.
Crucially, $\alpha$ depends \emph{only on~$W$}, not on~$p$.

\emph{Step 2: Second derivative in $z$.} Differentiate
\eqref{eq:thm-h-v2} a second time in $z$ at $z=0$. The second
derivative of the left-hand side involves $h''(W)$ applied to the
squares of the respective chain-rule factors:
\[
   h''(W)\cdot\!\left[\!\left(\frac{1-p}{\alpha(W)}\right)^{\!2}
   - \left(\frac{p}{\alpha(W)}\right)^{\!2}\right]
   \;=\; h''(W)\cdot\frac{(1-p)^{2}-p^{2}}{\alpha(W)^{2}}
   \;=\; h''(W)\cdot\frac{1-2p}{\alpha(W)^{2}}.
\]
The right-hand side of \eqref{eq:thm-h-v2} is $-z$, which is linear in
$z$; its second derivative is zero. Therefore
\[
   h''(W)\cdot\frac{1-2p}{\alpha(W)^{2}} \;=\; 0
   \qquad \text{for all } p\in(0,1).
\]
Choosing any $p\neq\tfrac12$ forces $h''(W)=0$ for every
$W\in\mathcal{D}$.

\emph{Step 3: Integration.} Since $h''=0$ on $\mathcal{D}$, the function
$h$ is affine: $h(W)=-\alpha_{0}W+C$ for constants $\alpha_{0}>0$ (from
$\alpha=-h'>0$) and $C\in\R$.
Hence $U'(W)=e^{h(W)}=e^{C}\exp(-\alpha_{0}W)$. Integrating:
\[
   U(W) \;=\; -\frac{e^{C}}{\alpha_{0}}\,\exp(-\alpha_{0}W) + D
   \;=\; a - b\exp(-\alpha_{0}W),
\]
with $a\equiv D\in\R$ and $b\equiv e^{C}/\alpha_{0}>0$. This is the CARA
utility function, and $\alpha(p,W)\equiv\alpha_{0}$ is constant.

\emph{Step 4: Converse.} Set $U(W)=a-b\exp(-\alpha_{0}W)$ in
\eqref{eq:thm-foc-v2}. Then
$U'(W)=b\alpha_{0}\exp(-\alpha_{0}W)$, and the FOC ratio reduces to
\[
   \ln\frac{U'(W+(1-p)x)}{U'(W-px)}
   \;=\; -\alpha_{0}\bigl[(W+(1-p)x)-(W-px)\bigr]
   \;=\; -\alpha_{0}x,
\]
so $-\alpha_{0}x = -z$, giving $x=z/\alpha_{0}$, which is
\eqref{eq:linear-demand-rep} with $\alpha\equiv\alpha_{0}$.
\end{proof}

\begin{proof}[Proof of Proposition~\ref{prop:cara}]
Substituting the CARA demand $x_{k}=[\logit\mu_{k}-\logit p]/\alpha_{k}$
into the market-clearing condition $\sum_{k=1}^{K}x_{k}=0$:
\[
   \sum_{k=1}^{K} \frac{\logit\mu_{k}-\logit p}{\alpha_{k}} \;=\; 0.
\]
Solving for $\logit p$:
\[
   \logit p \cdot \sum_{k}\alpha_{k}^{-1}
   \;=\; \sum_{k}\alpha_{k}^{-1}\logit\mu_{k},
\]
hence $\logit p=\sum_{k}w_{k}\logit\mu_{k}$ with
$w_{k}=\alpha_{k}^{-1}/\sum_{j}\alpha_{j}^{-1}$.
By \eqref{eq:posterior}, $\logit\mu_{k}=\tau_{k}u_{k}$, so
\[
   \logit p \;=\; \sum_{k}w_{k}\tau_{k}u_{k}.
\]
This is a weighted linear combination of the signals. When
$\alpha_{k}\equiv\alpha$ for all $k$, $w_{k}=1/K$ and this reduces to
$\logit p = \Tstar/K$ where $\Tstar=\sum_{k}\tau_{k}u_{k}$. The price is
then informationally equivalent to $\Tstar$: any agent observing $p$
can infer $\Tstar$ by inverting $\Lam(\cdot/K)$. Since $\Tstar$ is
sufficient for $v$, the equilibrium is fully revealing.

With heterogeneous $\alpha_{k}$, the no-learning price reveals
$\sum_{k}w_{k}\tau_{k}u_{k}$, which for $K\geq 3$ is generically not
informationally equivalent to $\Tstar$ (different linear combinations
of $K\geq 3$ variables are linearly independent). However, the price
function $p=\Lam(\Tstar/K)$ is still a valid REE: at this price every
agent's posterior equals $\Lam(\Tstar/K)=p$, demands are zero
identically, and the market clears with zero supply. This is an REE
for any $(\alpha_{k})$.

The exponential-family extension replaces $\tau_{k}u_{k}$ by the
natural-parameter log-likelihood ratio of the chosen family; the
identical linear-aggregation argument applies.
\end{proof}

\begin{proof}[Proof of Proposition~\ref{prop:non-cara}]
Under log utility ($\gamma=1$), the demand simplifies to
$x_{k}=W(\mu_{k}-p)/[p(1-p)]$ (set $R=e^{z}$ in \eqref{eq:crra-demand};
for $\gamma=1$, $R-1=e^{z}-1$ and the denominator is
$(1-p)+e^{z}p$; direct computation yields the stated form).
Substituting into $\sum_{k}x_{k}=0$ with homogeneous $W$:
\[
   \sum_{k=1}^{K}\frac{\mu_{k}-p}{p(1-p)} \;=\; 0
   \qquad\Longrightarrow\qquad
   p \;=\; \bar\mu \;\equiv\; \frac{1}{K}\sum_{k}\mu_{k}
   \;=\; \frac{1}{K}\sum_{k}\Lam(\tau u_{k}).
\]
To show this is not a function of $\Tstar$ alone, fix $\Tstar=0$ and
consider two signal realisations with $K=3$:
\begin{itemize}[topsep=2pt,itemsep=2pt]
\item Realisation~A: $(u_{1},u_{2},u_{3})=(0,0,0)$. Then
$p_{A}=\Lam(0)=\tfrac12$.
\item Realisation~B: $(u_{1},u_{2},u_{3})=(\delta,-\delta,0)$ for any
$\delta\neq 0$. Then
$p_{B}=\tfrac13[\Lam(\tau\delta)+\Lam(-\tau\delta)+\Lam(0)]$.
\end{itemize}
Both have $\Tstar=\tau(0+0+0)=0=\tau(\delta-\delta+0)$, but $p_{A}\neq
p_{B}$ because
\[
   \tfrac12\bigl[\Lam(\tau\delta)+\Lam(-\tau\delta)\bigr]
   \;\neq\; \Lam(0) = \tfrac12
\]
by Jensen's inequality applied to the strictly concave--convex logistic
function: $\Lam$ is strictly concave on $(0,\infty)$ and strictly convex
on $(-\infty,0)$, so $\Lam(\tau\delta)+\Lam(-\tau\delta)<2\Lam(0)$ for
$\delta\neq 0$. Hence $p_{B}<p_{A}$, and $p$ depends on the empirical
distribution of $u_{k}$, not just on $\Tstar$.
\end{proof}

\begin{proof}[Proof of Proposition~\ref{prop:jensen}]
Taylor-expand the logistic function around zero to fifth order. Using
$\Lam'(0)=\tfrac14$, $\Lam''(0)=0$, $\Lam'''(0)=-\tfrac18$:
\[
   \Lam(x) \;=\; \tfrac12+\tfrac14 x - \tfrac{1}{48}x^{3}+O(x^{5}).
\]
Substituting $x=\tau u_{k}$ and averaging over $k=1,\ldots,K$:
\[
   p \;=\; \frac{1}{K}\sum_{k}\Lam(\tau u_{k})
   \;=\; \tfrac12+\tfrac{\tau}{4}\bar u
         -\frac{\tau^{3}}{48K}\sum_{k}u_{k}^{3}+O(\tau^{5}),
\]
where $\bar u=\tfrac{1}{K}\sum_{k}u_{k}$.
The CARA price at the same $\Tstar=\tau K\bar u$ is
\[
   p^{\mathrm{CARA}}
   \;=\; \Lam\!\left(\frac{\Tstar}{K}\right)
   \;=\; \Lam(\tau\bar u)
   \;=\; \tfrac12+\tfrac{\tau}{4}\bar u - \tfrac{\tau^{3}}{48}\bar u^{3}
         +O(\tau^{5}).
\]
Subtracting:
\[
   \Delta p \;=\; p - p^{\mathrm{CARA}}
   \;=\; -\frac{\tau^{3}}{48}\left[
         \frac{1}{K}\sum_{k}u_{k}^{3} - \bar u^{3}\right]
   +O(\tau^{5}).
\]
Using $\bar u=U_{1}/K$ where $U_{1}=\sum_{k}u_{k}$ and
$U_{3}=\sum_{k}u_{k}^{3}$, we obtain $\bar u^{3}=U_{1}^{3}/K^{3}$ and
$\tfrac{1}{K}\sum u_{k}^{3}=U_{3}/K$, giving
\eqref{eq:jensen}.
\end{proof}

\begin{proof}[Proof of Proposition~\ref{prop:smooth}]
\emph{Continuity.} For each $\gamma\in(0,\infty]$ the no-learning
equilibrium price $p_{\gamma}(u_{1},\ldots,u_{K})$ is the unique solution
of the market-clearing equation
$\sum_{k}x_{k}^{\gamma}(\Lam(\tau u_{k}),p)=0$, where $x_{k}^{\gamma}$
denotes the CRRA demand of \eqref{eq:crra-demand} at parameter $\gamma$.

For any fixed $(u_{1},\ldots,u_{K})$ and $p\in(0,1)$, the demand
$x_{k}^{\gamma}(\mu_{k},p)$ is jointly continuous in $(\gamma,p)$ on
$(0,\infty]\times(0,1)$. Moreover, $x_{k}^{\gamma}$ is strictly
decreasing in $p$ (buying less when the asset is more expensive). By the
implicit function theorem applied to the market-clearing equation
$G(\gamma,p)=\sum_{k}x_{k}^{\gamma}(\mu_{k},p)=0$, with
$\partial G/\partial p<0$, the solution $p_{\gamma}$ is a continuous
function of $\gamma$ on $(0,\infty]$ uniformly over compact subsets of
the signal support.

The CARA limit follows from Lemma~\ref{lem:cara-limit}:
$\lim_{\gamma\to\infty}x_{k}^{\gamma}/\gamma
=[\logit\Lam(\tau u_{k})-\logit p]/W_{k}$, so the normalised demands
converge to CARA demands and the normalised price converges to the CARA
price.

The continuity of $1-R^{2}(\Tstar;p_{\gamma})$ in $\gamma$ on
$(0,\infty]$ follows by dominated convergence: the random variables
$(\Tstar, \logit p_{\gamma})$ converge jointly in distribution as
$\gamma\to\infty$, and their second moments are uniformly bounded on a
compact signal support. The $\gamma=\infty$ value is $0$ by
Proposition~\ref{prop:cara}.

\emph{Strict positivity.} By Proposition~\ref{prop:non-cara}, the
log-utility ($\gamma=1$) equilibrium satisfies $1-R^{2}>0$. The
continuity just established implies that the deficit remains strictly
positive in a neighbourhood of $\gamma=1$.

More generally, the Jensen-gap expansion of
Proposition~\ref{prop:jensen} gives
$\Delta p=O(\tau^{3}/\gamma^{2})$ at small $\tau$ (the CRRA demand
approaches CARA as $\gamma\to\infty$, and the deviation from CARA
pricing scales as $\gamma^{-2}$ through the second-order term in the
Taylor expansion of $R=e^{z/\gamma}$). This is strictly positive for
every finite $\gamma$.

Since $1-R^{2}$ is a continuous function of $\Delta p$ (and vanishes
only when $\Delta p$ vanishes identically over the joint signal
distribution), the deficit is strictly positive for all
$\gamma<\infty$.

The deficit is also monotonically decreasing in $\gamma$ at every
$(\tau,\gamma)$ pair tested numerically (Table~\ref{tab:smooth}).
The monotonicity at small $\tau$ follows analytically from the
$O(\tau^{3}/\gamma^{2})$ scaling; the global claim (all $\tau$) rests on
the numerical evidence.
\end{proof}

\begin{proof}[Proof of Proposition~\ref{prop:ree}]
The proof has two parts: an analytical argument that the limit of the
Picard sequence is partially revealing if it exists, and a numerical
verification that the sequence converges.

\emph{Part (a): the limit is partially revealing.} Define the Picard
sequence $\mu^{0}(u,p)=\Lam(\tau u)$ (the no-learning posterior) and
$\mu^{n+1}=\Phi(\mu^{n})$, where $\Phi$ is the contour-integration
Bayes update defined in Section~\ref{sec:ree}. We show by induction
that at every $n\geq 0$, the demand function
$d^{n}(u;p)\equiv x(\mu^{n}(u,p),p)$ is strictly nonlinear in $u$ at
every interior price $p$, and hence the contour
$\{(u_{2},u_{3}):d^{n}(u_{2};p)+d^{n}(u_{3};p)=-d^{n}(u_{1};p)\}$ is
curved.

\emph{Base case} ($n=0$). Since $\mu^{0}=\Lam(\tau u)$, the demand is
$d^{0}(u)=x(\Lam(\tau u),p)$. Write $z=\logit(\Lam(\tau u))-\logit(p)
=\tau u-\logit p$. The CRRA demand as a function of $z$ is
\[
   x(z) \;=\; \frac{W(e^{z/\gamma}-1)}{(1-p)+e^{z/\gamma}p}.
\]
Computing the second derivative: let $R=e^{z/\gamma}$, then
\begin{align*}
   \frac{\partial x}{\partial z}
   &\;=\; \frac{W}{\gamma}\cdot\frac{R}{[(1-p)+Rp]^{2}}, \\
   \frac{\partial^{2} x}{\partial z^{2}}
   &\;=\; \frac{W}{\gamma^{2}}\cdot\frac{R\,(1-p-Rp)}{[(1-p)+Rp]^{3}}.
\end{align*}
The numerator of $\partial^{2}x/\partial z^{2}$ is $R(1-p-Rp)$, which
is zero only at $R=(1-p)/p$, i.e.\ at $z=-\gamma\logit p$. At all other
values of $z$, the second derivative is nonzero. Since $z=\tau u-\logit p$
is affine in $u$, the composition $d^{0}(u)=x(\tau u-\logit p)$ inherits
this nonlinearity: $\partial^{2}d^{0}/\partial u^{2}\neq 0$ at all but
one value of $u$.

Under CARA, $x(z)=z/\alpha$ is affine in $z$, so $d^{0}$ is affine in
$u$ and the contour is straight---that is the content of
Proposition~\ref{prop:cara}.

\emph{Inductive step.} Suppose $d^{n}$ is nonlinear in $u$. Then:
\begin{enumerate}[label=(\roman*),topsep=2pt,itemsep=2pt]
\item The contour at $(u_{1},p)$ is curved, so $u_{2}+u_{3}$ varies along
it. The density ratio
\[
   \frac{f_{1}(u_{2})\,f_{1}(u_{3})}{f_{0}(u_{2})\,f_{0}(u_{3})}
   \;=\; \exp\!\bigl(\tau(u_{2}+u_{3})\bigr)
\]
is therefore non-constant on the contour. (Under CARA, $u_{2}+u_{3}$ is
constant on the straight contour, so the ratio is constant and the
agent learns $\Tstar$ exactly.)

\item The contour integrals
$A_{v}(u_{1},p)=\int_{C}f_{v}(u_{2})\,f_{v}(u_{3}^{\star}(u_{2}))\,du_{2}$
average the non-constant density ratio against the signal density. The
log-ratio $\log(A_{1}/A_{0})$ depends on $u_{1}$ because different own
signals $u_{1}$ generate different contour shapes (the agent-1 slice
$P[u_{1},\cdot,\cdot]$ varies with $u_{1}$). The updated posterior
\[
   \mu^{n+1}(u_{1},p) \;=\; \Lam\!\bigl(\tau u_{1}+\log(A_{1}/A_{0})\bigr)
\]
thus depends on $u_{1}$ beyond the price---it is not identically equal
to $p$.

\item The demand $d^{n+1}(u)=x(\mu^{n+1}(u,p),p)$ is the composition of
the CRRA demand (which has nonzero second derivative in $z$ for finite
$\gamma$ at all but one point, as shown in the base case) with
$\mu^{n+1}(u,p)$ (which depends nontrivially on $u$ by (ii)). For this
composition to be affine in $u$, the curvature of $\mu^{n+1}$ would need
to exactly cancel the demand curvature at the inflection point
$z=-\gamma\logit p$ for every $u$ simultaneously---a codimension-one
condition that is violated at every iterate and grid resolution tested
numerically.
\end{enumerate}

Since $d^{n}$ is nonlinear at every $n$, the contour is curved at every
$n$. If $\mu^{n}\to\mu^{\star}$, then $d^{n}\to d^{\star}$, and the
nonlinearity bound passes to the limit by continuity of the
second-derivative test. Hence $d^{\star}$ is nonlinear, the limit
contour is curved, and the limit equilibrium is partially revealing.

\emph{Part (b): convergence.} The Picard sequence is iterated
numerically using the posterior-function method of
Appendix~\ref{app:numerics} with isotonic-regression monotonicity
projection and Newton--Krylov polishing. At grid resolution $G=20$
with 50-digit multiprecision arithmetic, convergence
$\|F\|_{\infty}<10^{-25}$ is achieved. The converged deficit $1-R^{2}$
(measured via probability-weighted regression) is stable under grid
refinement from $G=10$ to $G=20$
(Table~\ref{tab:G-ladder}).
\end{proof}

\begin{proof}[Proof of Proposition~\ref{prop:welfare}]
\emph{No-learning equilibrium.} Under CARA at the no-learning
equilibrium of Proposition~\ref{prop:cara} with homogeneous parameters,
$\logit\mu_{k}=\tau u_{k}$ and $\logit p=(1/K)\sum_{k}\tau u_{k}$, so
\[
   \logit\mu_{k}-\logit p
   \;=\; \tau u_{k} - \frac{1}{K}\sum_{j}\tau u_{j}
   \;=\; \tau(u_{k}-\bar u).
\]
This is zero only when $u_{k}=\bar u$ for all $k$, which occurs on the
measure-zero set $\{u_{1}=\cdots=u_{K}\}$. The CARA demand
$x_{k}=\tau(u_{k}-\bar u)/\alpha$ is thus nonzero almost surely, and
since market clearing $\sum_{k}x_{k}=0$ forces both positive and
negative demands, aggregate volume
$V_{\mathrm{vol}}=\tfrac12\sum_{k}|x_{k}|$ is strictly positive on a
full-measure set.

Under CRRA, the private posteriors $\mu_{k}=\Lam(\tau u_{k})$ differ
generically across $k$. Market clearing pins the price in the interval
$p\in(\min_{k}\mu_{k},\max_{k}\mu_{k})$, so $\mu_{k}-p\neq 0$ for at
least one $k$ (in fact for all but at most one $k$) on a full-measure set.
The CRRA demand \eqref{eq:crra-demand} is nonzero whenever $\mu_{k}\neq p$
(since $R\neq 1$ when $z\neq 0$), so volume is strictly positive.

\emph{REE.} At the CARA REE, all posteriors coincide with the
fully-revealing price: $\mu_{k}=p=\Lam(\Tstar/K)$ for every $k$ and
every signal realisation. Therefore
$z_{k}=\logit\mu_{k}-\logit p=0$ for all $k$, and the CARA demand
is $x_{k}=0/\alpha=0$ identically. This is the \citet{Milgrom1982}
no-trade outcome: when the price is fully revealing, there are no gains
from trade.

At the CRRA REE of Proposition~\ref{prop:ree}, posteriors disagree on
a full-measure set (the converged $\mu_{k}^{\star}(u_{k},p)$ differ
across agents because the contour integrals differ). Since
$\mu_{k}\neq p$ for some $k$, the CRRA demand is nonzero and aggregate
volume is strictly positive.
\end{proof}

\begin{proof}[Proof of Proposition~\ref{prop:value}]
The certainty-equivalent value of acquiring an additional
precision-$\tau$ signal $u_{0}$ at the converged REE is
\[
   V(\tau) \;=\;
   \E\!\left[\,\mathrm{CE}\!\left(\mu^{+}\!,p\right)
            -\mathrm{CE}\!\left(\mu^{-}\!,p\right)\,\right],
\]
where $\mu^{+}$ is the posterior of an agent who additionally observes
$u_{0}$, $\mu^{-}$ is the posterior of an agent who does not, $p$ is the
equilibrium price, and $\mathrm{CE}(\mu,p)$ denotes the certainty
equivalent of optimal trading at posterior $\mu$ and price $p$.

\emph{Under CARA.} At the REE, $p=\Lam(\Tstar/K)$ already encodes the
joint posterior. The price is a sufficient statistic for $v$ given the
existing signals. Therefore $\mu^{+}=\mu^{-}=\Lam(\Tstar/K)$ for
almost every realisation that conditions on the price: the additional
signal $u_{0}$ is redundant. The integrand
$\mathrm{CE}(\mu^{+},p)-\mathrm{CE}(\mu^{-},p)=0$ almost surely, so
$V_{\mathrm{CARA}}(\tau)=0$ for every $\tau$.

\emph{Under CRRA.} At the REE, $p$ encodes $\Tstar$ only partially
(Proposition~\ref{prop:ree}). The additional signal $u_{0}$ provides
information about $v$ that is not already in the price, so
$\mu^{+}\neq\mu^{-}$ on a positive-measure set.

Non-negativity: An agent can always replicate
$\mathrm{CE}(\mu^{-},p)$ by ignoring the additional signal (choosing
the same trade she would without it), so $\mathrm{CE}(\mu^{+},p)\geq
\mathrm{CE}(\mu^{-},p)$ and the integrand is non-negative everywhere.

Strict positivity: On the positive-measure set where $\mu^{+}\neq
\mu^{-}$, the agent can improve her trade by conditioning on the
additional signal. The optimal demand $x^{+}\neq x^{-}$ (since the FOC
at $\mu^{+}\neq\mu^{-}$ has a different solution), and the certainty
equivalent is strictly higher at the optimal trade than at the
suboptimal one. Hence
$\mathrm{CE}(\mu^{+},p)>\mathrm{CE}(\mu^{-},p)$ on this set, and
$V_{\mathrm{CRRA}}(\tau)>0$.

For the small-$\tau$ derivative: a Taylor expansion of the posterior in
$\tau$ gives $\mu^{\pm}-\bar\mu=O(\tau)$ with leading-order coefficient
proportional to $u_{0}$ for $\mu^{+}$ and $0$ for $\mu^{-}$. The
certainty-equivalent shift is order $\tau^{2}$ (from the second-order
condition for optimal trading), so
$V_{\mathrm{CRRA}}'(0^{+})>0$.
\end{proof}

\begin{proof}[Proof of Proposition~\ref{prop:gs}]
\emph{Non-existence under CARA.} Suppose for contradiction that there
exists a CARA REE with no noise and a strictly positive measure
$\lambda^{\star}\in(0,1]$ of agents acquiring the signal at cost $c>0$.
By Proposition~\ref{prop:cara} and the homogeneity of the acquiring
group, the price function in such an equilibrium is fully revealing:
$p=\Lam(\Tstar/K)$ over the acquiring group. By
Proposition~\ref{prop:value}, $V_{\mathrm{CARA}}(\tau)=0$ at this fully
revealing price, so each acquirer's gross value of information is zero.
The net payoff from acquisition is $0-c=-c<0$. Every acquirer strictly
prefers not to acquire. Hence $\lambda^{\star}>0$ cannot be an
equilibrium---contradiction.

If $\lambda=0$ (no one acquires), the price is uninformative and
$V_{\mathrm{CARA}}(\tau,0)>0$, so each agent strictly prefers to acquire.
Hence $\lambda=0$ is not an equilibrium either. No equilibrium exists:
this is the \citet{GS1980} paradox.

\emph{Existence under CRRA.} Let $\lambda\in[0,1]$ denote the fraction of
agents acquiring the signal, and let $V_{\mathrm{CRRA}}(\tau,\lambda)$
denote the marginal value of acquisition at fraction $\lambda$.

At $\lambda=0^{+}$: the price reveals nothing, so the first acquirer
reaps the full informational rent. By Proposition~\ref{prop:value},
$V_{\mathrm{CRRA}}(\tau,0^{+})>0$.

At $\lambda=1$: all agents have the signal. The price is partially
revealing (Proposition~\ref{prop:ree}), not fully revealing. The
marginal value $\bar c\equiv V_{\mathrm{CRRA}}(\tau,1)$ is strictly
positive but smaller than $V_{\mathrm{CRRA}}(\tau,0^{+})$ (more
acquirers $\Rightarrow$ more informative price $\Rightarrow$ smaller
marginal informational rent).

$V_{\mathrm{CRRA}}(\tau,\lambda)$ is continuous in $\lambda$ (by
continuity of the equilibrium price function in $\lambda$) and weakly
decreasing (each additional acquirer makes the price weakly more
informative). For any acquisition cost $c\in(0,\bar c)$, the equation
$V_{\mathrm{CRRA}}(\tau,\lambda)=c$ has at least one solution
$\lambda^{\star}(c)\in(0,1)$ by the intermediate-value theorem. At this
$\lambda^{\star}$, the marginal acquirer is exactly indifferent between
acquiring and not acquiring, and no agent has an incentive to deviate.
This is a Grossman--Stiglitz equilibrium with positive information
acquisition and partially revealing prices.
\end{proof}
